\thispagestyle{empty}

\begin{center}
  {\bf \Huge Ringraziamenti}
\end{center}

\vspace{4cm}
\emph{
  Chi mi conosce sa che non sono mai stato di molte parole, 
  ma desidero comunque ringraziare tutte le persone che mi hanno supportato ed aiutato 
  nel percorso che ha portato a questa tesi: dal Professor Bouquet, che 
  mi ha guidato nella stesura di un testo che raccontasse il mio lavoro,
  a chi, come la mia famiglia, mi ha sopportato mentre cercavo di spiegare
  quale fosse l'argomento affrontato. 
  È poi dovuto un grazie anche a tutti i compagni di corso che durante questi
  tre anni sono diventati amici, rendendo piacevole anche la frequentazione
  dei corsi peggiori, oltre che alla mia bella ragazza, che con la sua compagnia 
  è sempre stata capace di convincermi ad andare a studiare.
  Ultimi ma non per importanza, ringrazio tutti i colleghi di FBK (da cui è nato questo lavoro),
  tutti quelli che mi hanno fatto appassionare a questi temi e comunque chiunque si senta di
  poter essere ringraziato: le persone sono tante, ma c'è un pensiero per tutte.
}


\vspace*{\fill}
\begin{flushright}
\begin{minipage}{0.58\textwidth}
  \raggedleft\itshape\small
  Non si può capire un processo arrestandolo.
  La comprensione deve fluire insieme col processo,
  deve unirsi a esso e fluire con esso.
  \par\vspace{0.9em}
  \upshape\footnotesize
  Prima Legge del Mentat\par
  \textsc{Frank Herbert}, \emph{Dune}\par
\end{minipage}
\end{flushright}
\cleardoublepage
